% !TeX root = ../main.tex
\chapter{Literature Review}
\label{ch:literature}
\section{Training Strategy and Deployment Objective}
The methodological background is organised around the project's separation of
training and deployment. A virtual vehicle supplies interaction, RL learns a
teacher, imitation creates a sensor-limited student, and a compact model provides
the candidate for board inference. The objective is to replace manually supplied
driving examples while retaining useful behaviour in a constrained implementation.
The subsequent model comparison evaluates this training strategy rather than
introducing a new tree or policy-gradient algorithm.

Kritzinger et al. distinguish digital models, digital shadows and digital twins
according to their level of integration~\cite{kritzinger2018digital}. This
distinction prevents a simulator from being treated as evidence of a fully
connected physical--digital system. In this dissertation, the virtual vehicle
implements the training side of an intended digital-twin workflow; measured
physical calibration and synchronisation are still required.

The digital-twin-oriented scope is operationally defined in
Chapter~\ref{ch:introduction}: a vehicle-specific virtual counterpart supports
learning and collection, while calibration and physical coupling remain pending.
The cited simulation platform supplies the enabling environment~\cite{koenig2004gazebo};
its existence alone does not establish physical fidelity for this vehicle.
Likewise, privileged-policy transfer motivates the separation of teacher and
student information~\cite{chen2019cheating}, but provides no substitute for
measuring the present student's behaviour.

\section{Local Feedback Control}
The engineering question in this study concerns a local feedback controller:
a range observation is converted into a steering decision, which changes the
next observation. This differs from evaluating a classifier on an independent
collection of labelled samples. It also differs from global navigation, where
route choice and destination progress require additional task definitions.
The literature considered here is selected to explain teacher learning,
student imitation, tree-based inference and simulation perturbations.

A classical example is the Vector Field Histogram (VFH) of Borenstein and
Koren~\cite{borenstein1991vfh}. It maintains a Cartesian obstacle histogram
and derives a polar representation from which a steering direction is selected.
The present student instead maps a single binned scan directly to an action
learned from demonstrations. Both approaches use directional range information,
but the student does not implement VFH's obstacle grid or selection rule.
VFH is methodological context here; no experimental comparison with VFH has
been conducted.

Tai, Paolo and Liu demonstrate a different route to obtaining navigation
behaviour without prior driving demonstrations: a deep RL policy maps ten range
measurements and a relative target position to continuous steering commands,
with reported tests in virtual and real environments~\cite{tai2017virtual}.
This supports investigating simulation interaction as a source of behaviour.
Their goal-directed neural controller differs from the present forward-cruising,
discrete-action task. Here the learned teacher supplies examples for a separate
tree student, and the reported evidence remains simulation-only. Their physical
transfer results cannot be attributed to the present vehicle.

For the present vehicle, the sensor representation is part of the control
design. Discretising distance by angle fixes the relationship between input
columns and physical directions. Suppressing the rear sector reduces available
information, while using zero for missing returns makes missingness a feature
of the learned policy. A comparison of learning algorithms is meaningful only
when these conventions and the action interface remain explicit.

\section{Random Forests and Compact Tree Policies}
Breiman introduced random forests as ensembles of randomised tree predictors
and analysed their generalisation behaviour in relation to the strength and
dependence of individual trees~\cite{breiman2001random}. The relevant design
idea is that multiple trees can provide different partitions of the input
space. An ensemble nevertheless requires more inference work than a single
tree, and its performance depends on how those partitions are constructed.

In this project, depth and tree count are explicit model choices. A shallow
forest permits several small decision structures to contribute to a prediction;
a deeper single tree permits a longer sequence of conditional decisions along
one path. These are different ways of allocating capacity. Neither allocation
can be declared superior for the vehicle without checking the states reached
during closed-loop operation.

The implementation uses scikit-learn 1.7.2, from the machine-learning library
described by Pedregosa et al.~\cite{pedregosa2011sklearn}. Its random-forest classifier
aggregates class probabilities from the trees, with bootstrap sampling and
random feature selection available during fitting~\cite{sklearnForest}.
This detail matters when exporting inference: replacing probability averaging
with hard majority voting would define a different controller. The study
therefore verifies agreement after export rather than treating any generated
C++ tree code as automatically equivalent.

Embedded tree learning also has more specialised alternatives. Kumar, Goyal
and Varma's Bonsai jointly learns a sparse projection and a shallow tree with
expressive nodes, targeting devices with as little as 2~KB of
RAM~\cite{kumar2017bonsai}. Its design shows why resource constraints should
affect model construction, rather than being checked only after fitting.
Bonsai is not an ordinary random forest, and its published memory results do
not describe the exported models in this project. The present forest provides
a simpler existing fitting/export path; the single-tree comparison tests
whether that ensemble complexity is justified by the observed behaviour.

Nor are trees the only feasible microcontroller models. Lai, Suda and Chandra
develop CMSIS-NN kernels to improve neural inference efficiency on Arm Cortex-M
CPUs~\cite{lai2018cmsis}. Thus a compact neural student is a credible alternative,
especially on a Cortex-M7 with floating-point support. Selecting a forest here
rests on its threshold representation, bounded traversal and verified C++ export,
not on an assertion that neural networks cannot execute on the target board.
An equal-budget neural comparison remains future work.

\section{Learning a Privileged Teacher}
PPO uses a policy-gradient objective with a clipped probability ratio to
moderate updates relative to the data-collecting policy~\cite{schulman2017ppo}.
Here it supplies the teacher policy rather than the final tree model. A neural
policy can use the simulated distance field and vehicle state while generating
examples for a more restricted student.

The separation of a privileged policy from a sensor-limited policy is related
to the two-stage approach in \emph{Learning by Cheating}. Chen et al. first
learn a policy with privileged information and then transfer behaviour to an
agent using ordinary observations in a driving simulator~\cite{chen2019cheating}.
The present project differs in its small differential-drive platform, LiDAR
bins, discrete steering labels and tree-based student. The cited work motivates
the information separation; it does not validate the present implementation.

A privileged teacher should not be interpreted as an available physical
controller. Its geometry is supplied by the simulator, and the student does
not acquire that geometry merely by imitating actions. Evaluating teacher and
student as if they had identical information would conceal a central design
constraint.

Rusu et al. show that an RL policy can be transferred to a smaller neural
network through policy distillation~\cite{rusu2015distillation}. Bastani, Pu and
Solar-Lezama's VIPER instead extracts a decision-tree policy from a neural
oracle using its Q-function, enabling verification in their benchmark
settings~\cite{bastani2018viper}. Together these works motivate separating the
model that learns the task from the model chosen for deployment. The present
pipeline uses hard action labels and ordinary supervised fitting; it does not
implement VIPER's Q-guided extraction or establish its formal guarantees.
The match to the research objective is the teacher--student separation, while
the collection procedure and validation claims remain specific to this project.

\section{Behaviour Cloning and Distribution Shift}
Ross, Gordon and Bagnell show why sequential prediction differs from ordinary
independent supervised prediction: actions influence future observations.
Their DAgger approach collects expert-labelled examples on states encountered
by the learner~\cite{ross2011reduction}. The distinction is relevant when
interpreting a student that performs well on teacher scans but fails after its
own steering changes the trajectory.

The present pipeline performs behaviour cloning on teacher-driven episodes.
Perturbed and near-wall teacher starts expand the demonstration conditions,
but there is no iterative learner-execution and expert-relabeling loop.
Consequently, the implementation is not described as DAgger. Likewise, labels
are discrete teacher commands, not teacher probability distributions; no
soft-label distillation loss is used.

This methodological boundary leads to a practical evaluation requirement.
Prediction disagreement is measured offline, but collisions and recovery
must be measured while the student actually drives. High agreement on common
straight-driving examples cannot establish adequate behaviour in a rare
turning state or after a loss of sensor coverage.

\section{Simulation Perturbations and Transfer}
Tobin et al. investigate domain randomisation by varying rendered simulation
conditions to support transfer of vision-based learning to physical
robotics~\cite{tobin2017domain}. That work establishes a motivation for varying
training conditions, but its image-based transfer result is not evidence that
any particular LiDAR corruption will improve the controller studied here.

The present experiments distinguish scattered point loss from contiguous
sector loss. These perturbations remove information in different spatial
patterns. A training procedure that masks a small number of separated
measurements need not provide examples of an extended blind region. The
experimental question is therefore whether the chosen augmentation matches
useful deployment variation, not simply whether augmented data are present.

Simulation also provides exact state and an idealised velocity interface.
Physical transfer would introduce sensor calibration, command-response and
timing questions. Those measurements are outside this draft, so simulation
robustness is discussed within the implemented perturbation families.

\section{Positioning of the Study}
Table~\ref{tab:literature} connects the selected methodological sources to the
choices made in this project. This is a focused initial review, not a claim
that these sources exhaust LiDAR navigation or embedded machine learning.
The contribution is the experimentally traceable combination of established
methods under a common observation, action and evaluation contract.

\begin{table}[htbp]
\centering\small
\caption{Methodological context and its role in this study.}
\label{tab:literature}
\begin{tabularx}{\textwidth}{p{3.1cm}>{\raggedright\arraybackslash}X>{\raggedright\arraybackslash}X}
\toprule
Source & Relevant idea & Boundary in this study \\
\midrule
Breiman~\cite{breiman2001random} & Randomised tree ensembles & Forest superiority is tested, not assumed. \\
Schulman et al.~\cite{schulman2017ppo} & Clipped policy optimisation & PPO trains the privileged teacher. \\
Chen et al.~\cite{chen2019cheating} & Privileged-to-sensor policy transfer & Student and task differ from the cited driving system. \\
Ross et al.~\cite{ross2011reduction} & Learner-induced state distribution & No iterative DAgger collection is implemented. \\
Tobin et al.~\cite{tobin2017domain} & Training-condition variation & LiDAR masking is assessed separately from physical transfer. \\
\bottomrule
\end{tabularx}
\end{table}

The reviewed work leads to a coherent selection: simulation supplies repeated
interaction; PPO provides the neural teacher; behaviour cloning transfers
actions to the LiDAR interface; and bounded-depth trees expose a manageable
inference structure. Each selection addresses a different constraint. The
literature motivates this combination but does not establish that it is the
best-performing or lowest-power combination for the vehicle.

The empirical gap addressed by the experiments is narrow and concrete:
whether automatically generated teacher demonstrations produce useful compact
students, and whether the initially selected forest retains its appeal when
offline agreement and closed-loop behaviour are assessed separately. A broader survey of
microcontroller navigation systems and directly comparable physical platforms
remains appropriate for the final dissertation, when the hardware scope and
measurements are settled.
